\documentclass[10pt]{article}
\usepackage[paperwidth=42cm, paperheight=30cm, margin=2cm]{geometry} % A3 Landscape
\usepackage{fontspec}
\setmainfont{Charter}
\newfontfamily\headerfont{Helvetica Neue}
\usepackage[brazil]{babel}

\usepackage{graphicx}
\usepackage{xcolor}
\usepackage{amsmath}
\usepackage{amssymb}
\usepackage{tikz}
\usetikzlibrary{calc, positioning, arrows.meta, shapes.geometric, decorations.pathreplacing, fit, backgrounds}

% --- PALETA DE CORES SÓBRIA (Newspaper Style) ---
\definecolor{navy}{RGB}{15, 23, 42}       % Texto principal e acento
\definecolor{accent}{RGB}{2, 132, 199}     % Azul para saturação (único ponto de cor)
\definecolor{lightgrey}{RGB}{241, 245, 249} % Fundo sutil
\definecolor{border}{RGB}{148, 163, 184}   % Linhas de separação
\definecolor{redaccent}{RGB}{153, 27, 27}  % Vermelho sóbrio para o LED

% Custom Bullet
\newcommand{\titem}[1]{\hangindent=2.5ex\raisebox{0.2ex}{\textcolor{navy}{\footnotesize$\blacktriangleright$}}\hspace{0.8ex}#1\par\vspace{0.6ex}}

\pagestyle{empty}

% MACROS DE CAIXAS ESTILO COLUNA DE JORNAL
\newcommand{\newsbox}[2]{
    \noindent
    \begin{minipage}{\linewidth}
        \vspace{0.4cm}
        {\headerfont\bfseries\large\textcolor{navy}{\uppercase{#1}}}\\[0.5ex]
        \drawline\\[1.5ex]
        {\small #2}
        \vspace{0.4cm}
    \end{minipage}\par
}

\newcommand{\drawline}{\noindent\textcolor{border}{\rule{\linewidth}{0.6pt}}}

\begin{document}
\begin{center}
\vspace*{\fill}
\scalebox{0.68}{%
\begin{minipage}{\textwidth}

% =========================================================
% TOP: HEADER ESTILO JORNAL
% =========================================================
\begin{center}
    {\headerfont\fontsize{54}{60}\bfseries\selectfont \uppercase{Oximetria de Pulso}}\\[0.4cm]
    \textcolor{border}{\rule{18cm}{2pt}}\\[0.3cm]
    {\large\textit{Protocolologia Crítica e Monitorização Técnica \textbullet\ por Nick Mark MD \textbullet\ Tradução Profissional}}\\[0.8cm]
\end{center}

\noindent
\begin{minipage}[t]{0.7\textwidth}
    \headerfont\textbf{\large PRINCÍPIO FISIOLÓGICO:}\\
    {\small A oximetria de pulso constitui a monitorização contínua da saturação de oxigênio da hemoglobina através de espectrofotometria. O dispositivo analisa a absorbância diferencial entre a \textbf{hemoglobina oxigenada (OxyHb)} e a \textbf{desoxigenada (DeoxyHb)} ao submeter o tecido a comprimentos de onda específicos nos espectros \textbf{vermelho} e \textbf{infravermelho}.}
\end{minipage}\hfill
\begin{minipage}[t]{0.28\textwidth}
    \begin{flushright}
        \headerfont\textbf{\small INDEX: PALS-ICU-024}\\
        \drawline\\
        {\footnotesize onepagericu.com \textbullet\ @nickmmark}
    \end{flushright}
\end{minipage}

\vspace{1.2cm}

% =========================================================
% DIAGRAMAS TÉCNICOS (DEDO E GRÁFICOS AGRUPADOS)
% =========================================================
\noindent
\begin{tikzpicture}[remember picture, overlay, shift={(current page.center)}]
    % Placeholder for layout matching, logic goes inside minipages below for flow
\end{tikzpicture}

\noindent
\begin{minipage}[t]{0.5\textwidth}
    \newsbox{Mecânica do Sensor}{
        \begin{minipage}[c]{0.45\linewidth}
            % Redesenho técnico do dedo
            \begin{tikzpicture}[scale=1.1]
                % Dedo (Contorno técnico)
                \draw[thick, navy!60, fill=lightgrey!50] (-2, -0.7) .. controls (1, -0.7) and (1.5, -0.5) .. (1.8, 0) .. controls (1.5, 0.5) and (1, 0.7) .. (-2, 0.7);
                % Unha
                \draw[navy!40] (0.8, 0.5) .. controls (1.5, 0.5) and (1.8, 0.3) .. (1.8, 0) .. controls (1.8, -0.3) and (1.5, -0.5) .. (0.8, -0.5);
                
                % Sensor superior
                \draw[thick, navy] (-0.5, 1.2) -- (0.8, 1.2) -- (0.8, 0.8) -- (-0.5, 0.8) -- cycle;
                \node at (0.15, 1.0) {\tiny\headerfont LED};
                
                % Sensor inferior
                \draw[thick, navy] (-0.5, -1.2) -- (0.8, -1.2) -- (0.8, -0.8) -- (-0.5, -0.8) -- cycle;
                \node at (0.15, -1.0) {\tiny\headerfont SENSOR};
                
                % Raios de luz
                \draw[dashed, redaccent, -Latex] (0.15, 0.8) -- (0.15, -0.8);
                \node[right, redaccent] at (0.15, 0.4) {\tiny 660nm};
                \draw[dashed, navy, -Latex] (0.45, 0.8) -- (0.45, -0.8);
                \node[right, navy] at (0.45, -0.4) {\tiny 940nm};
            \end{tikzpicture}
        \end{minipage}\hfill
        \begin{minipage}[c]{0.5\linewidth}
            Luz emitida atravessa o leito capilar pulsátil. O fotodetector oposto mede a intensidade da luz não absorvida.\\[1ex]
            \titem{O componente \textbf{AC} (pulsátil) isola o sangue arterial.}
            \titem{O componente \textbf{DC} (constante) filtra tecidos, ossos e sangue venoso.}
        \end{minipage}
    }
    
    \newsbox{Espectro de Absorção}{
        \centering
        \begin{tikzpicture}[scale=0.85, font=\headerfont]
            \draw[->, border] (0,0) -- (6,0) node[right] {\tiny nm};
            \draw[->, border] (0,0) -- (0,4) node[above] {\tiny Abs};
            
            % Curvas
            \draw[ultra thick, accent] (0, 3.5) .. controls (2, 2.8) and (3, 2) .. (5.5, 1) node[right] {\tiny OxyHb};
            \draw[ultra thick, redaccent!70] (0, 1) .. controls (1, 2) and (2, 3) .. (3, 3) .. controls (4, 3) and (5, 2.2) .. (5.5, 2.2) node[right] {\tiny DeoxyHb};
            
            % Linhas verticais técnicas
            \draw[dotted, navy] (0.8, 0) -- (0.8, 3.8) node[above, scale=0.7] {RED};
            \draw[dotted, navy] (4.8, 0) -- (4.8, 3.8) node[above, scale=0.7] {IR};
            
            \node[below, scale=0.6] at (0.8, 0) {660};
            \node[below, scale=0.6] at (4.8, 0) {940};
        \end{tikzpicture}
    }
\end{minipage}\hfill
\begin{minipage}[t]{0.45\textwidth}
    \newsbox{Forma de Onda (Pletismografia)}{
        \begin{minipage}[c]{\linewidth}
            \centering
            % Monitor clean
            \begin{tikzpicture}
                \draw[navy!40, fill=lightgrey!30, rounded corners=2pt] (0,0) rectangle (10, 4.5);
                
                % Waveflow
                \draw[thick, accent] (0.5, 1) .. controls (1.5, 1) and (1.8, 3) .. (2, 3) .. controls (2.2, 3) and (2.4, 2) .. (2.6, 2) .. controls (2.8, 2.2) and (3.1, 1) .. (4, 1)
                                      .. controls (5.0, 1) and (5.3, 3) .. (5.5, 3) .. controls (5.7, 3) and (5.9, 2) .. (6.1, 2) .. controls (6.3, 2.2) and (6.6, 1) .. (7.5, 1);
                
                \node[font=\headerfont\bfseries, accent] at (8.5, 3.5) {\Large 98};
                \node[font=\headerfont, scale=0.6, navy] at (8.5, 4) {SpO2 \%};
                
                \node[font=\headerfont\bfseries, navy!70] at (8.5, 1.5) {\Large 72};
                \node[font=\headerfont, scale=0.6, navy] at (8.5, 2) {PR bpm};
                
                % Annotations
                \draw[dashed, border] (2,3) -- (2,0.5);
                \node[below, scale=0.5, navy] at (2, 0.5) {SÍSTOLE};
                \draw[dashed, border] (2.6,2) -- (2.6,0.5);
                \node[below, scale=0.5, navy] at (2.6, 0.5) {NÓ DICRÓTICO};
            \end{tikzpicture}
        \end{minipage}\\[1.5ex]
        {\small A análise da onda pletismográfica valida a confiabilidade do valor numérico. Uma curva bem definida com nó dicrótico visível sugere perfusão adequada.}
    }
\end{minipage}

\vspace{0.8cm}
\drawline
\vspace{0.4cm}

% =========================================================
% TIER INFERIOR: COLUNAS DE JORNAL
% =========================================================
\noindent
\begin{minipage}[t]{0.31\textwidth}
    \newsbox{Fatores de Precisão}{
        \titem{\textbf{Pigmentação da Pele:} Melanina absorve luz. A SpO2 pode superestimar a saturação real em pacientes de pele negra, ocultando hipoxemias graves.}
        \titem{\textbf{Esmaltes:} Cores escuras (Preto, Azul, Verde) interferem significativamente na leitura.}
        \titem{\textbf{Perfusão:} Choque e vasoconstrição severa reduzem o sinal pulsátil.}
    }
\end{minipage}\hfill
\begin{minipage}[t]{0.32\textwidth}
    \newsbox{Interferências Químicas}{
        \titem{\textbf{Carboxi-Hb:} Falsamente normal. O sensor não distingue CO da O2.}
        \titem{\textbf{Met-Hb:} Leitura "travada" em 85\% independente do valor real.}
        \titem{\textbf{Corantes:} Azul de metileno causa queda súbita e falsa no SpO2.}
    }
\end{minipage}\hfill
\begin{minipage}[t]{0.31\textwidth}
    \newsbox{Limitações Técnicas}{
        \titem{\textbf{Atraso (Lag):} Periférico (5-15s). Durante ressuscitação, a leitura no monitor reflete o estado passado.}
        \titem{\textbf{Movimento:} Artefatos podem mimetizar ondas reais.}
        \titem{\textbf{Curva Shift:} Calibração limitada abaixo de 70\% de saturação.}
    }
\end{minipage}

\vspace{1.5cm}

\begin{minipage}[t]{0.45\textwidth}
    \headerfont\textbf{SUMÁRIO TÉCNICO:}\\
    {\small Sempre correlacione a forma de onda com o pulso palper. Em casos de dúvida sobre a precisão (Pele negra, intoxicação por CO, choque), a \textbf{gasometria arterial (SaO2)} permanece o padrão-ouro mandatório.}
\end{minipage}\hfill
\begin{minipage}[t]{0.45\textwidth}
    \begin{flushright}
        \textcolor{border}{\headerfont\fontsize{30}{30}\selectfont \uppercase{Critical Care Technical Series}}
    \end{flushright}
\end{minipage}

\end{minipage}%
}
\vspace*{\fill}
\end{center}
\end{document}
